%% \GlossRTL and \glscript under PDF/UA-2: the tags in reading order.
%%
%% The engines reverse the grid, never the macros, so the tree reads
%% column by column from the first word -- and a word's text must be its
%% LOGICAL text: XeTeX writes right-to-left glyphs in visual order, and
%% without /ActualText the first word would read "dliha" (finding 5 of
%% experiments/rtl; veraPDF passes that).  The Hebrew strings in the tree,
%% in order, are therefore exactly the script line and then the four words
%% as typed.  Languages: he on the script line and the four object words,
%% he-Latn on the four transliterated ones.
%% PDF/UA-2 preamble, mirroring examples/ua-demo.tex: this is the only
%% configuration whose output veraPDF is asked to validate, so it has to be
%% a realistic accessible build rather than a bare tagging switch.
%% pdfstandard=ua-2 requires PDF 2.0.  A dc:title is mandatory under
%% ISO 14289-2 clause 8.11.1, hence the pdftitle.
\DocumentMetadata{lang=en,pdfversion=2.0,pdfstandard=ua-2,tagging=on}
\documentclass{article}
\usepackage{iftex}
\ifPDFTeX
  \usepackage[T1]{fontenc}
  \usepackage[utf8]{inputenc}
\else
  \usepackage{fontspec}
\fi
\usepackage[lazy,gb4e]{linguexx}
\usepackage{hyperref}
\hypersetup{pdftitle={linguexx PDF/UA regression case},
            pdfdisplaydoctitle=true}
\pagestyle{empty}
\begin{document}

\newfontfamily\hebfont{FreeSerif.otf}[Script=Hebrew]
\GlossTierFont{1}{\hebfont}

\ex. \glscript{he}{\hebfont הילד אכל את התפוח}
\GlossTierLang{1}{he-Latn}\GlossTierFont{1}{\textnormal}
\gll ha-yeled axal et ha-tapuax \\
     \lpzg{def}-boy ate \lpzg{acc} \lpzg{def}-apple \\
\glt `The boy ate the apple.'

\GlossTierLang{1}{he}
\GlossTierLang{2}{he-Latn}
\GlossRTL
\ex. \glll הילד אכל את התפוח \\
           ha-yeled axal et ha-tapuax \\
           \lpzg{def}-boy ate \lpzg{acc} \lpzg{def}-apple \\
\glt `The boy ate the apple.'

\ex. NEXTEX stands where every example does.

\end{document}
